\documentclass{article}
\usepackage{amsmath,amssymb,xurl}
\usepackage[hidelinks]{hyperref}
% Shared commands for notes in docs/paper-gaps/.

\DeclareUnicodeCharacter{2080}{\ifmmode{}_{0}\else\textsubscript{0}\fi}
\DeclareUnicodeCharacter{2081}{\ifmmode{}_{1}\else\textsubscript{1}\fi}
\DeclareUnicodeCharacter{2082}{\ifmmode{}_{2}\else\textsubscript{2}\fi}
\DeclareUnicodeCharacter{2099}{\ifmmode{}_{n}\else\textsubscript{n}\fi}

\newcommand{\Lean}{\textsc{Lean}}
\ExplSyntaxOn
\NewDocumentCommand{\leanid}{m}{
  \ifmmode
    \textsf{\detokenize{#1}}
  \else
    \group_begin:
    \urlstyle{sf}
    \tl_set:Nn \l_tmpa_tl {#1}
    \tl_replace_all:Nnn \l_tmpa_tl {\_} {_}
    \exp_args:NV \path \l_tmpa_tl
    \group_end:
  \fi
}
\ExplSyntaxOff
\providecommand{\tr}{\operatorname{tr}}
\newcommand{\tnleanIssueBase}{https://github.com/LionSR/TNLean/issues/}
\newcommand{\tnleanPrBase}{https://github.com/LionSR/TNLean/pull/}
\newcommand{\tnleanDocBase}{https://sirui-lu.com/TNLean/docs/}
\newcommand{\ghissue}[1]{\href{\tnleanIssueBase#1}{GitHub issue \##1}}
\newcommand{\ghpr}[1]{\href{\tnleanPrBase#1}{PR \##1}}
\newcommand{\doclink}[1]{\href{\tnleanDocBase#1.html}{\path{#1}}}

% Dirac notation and the identity operator, as in blueprint/src/macros/common.tex.
% \providecommand keeps note-local definitions authoritative.
\usepackage{dsfont}
\providecommand{\ket}[1]{|#1\rangle}
\providecommand{\bra}[1]{\langle #1|}
\providecommand{\braket}[2]{\langle #1 | #2 \rangle}
\providecommand{\ketbra}[2]{|#1\rangle\!\langle #2|}
\providecommand{\Id}{\mathds{1}}
\providecommand{\one}{\mathds{1}}
\providecommand{\C}{\mathbb{C}}
\providecommand{\MN}[1]{M_{#1}(\C)}
\providecommand{\MD}{M_D(\C)}

% Verdict marker: \gapnote{<kind>}{<status>}. Kind is the policy
% classification (clarification, local-correction, scope-restriction,
% unfaithful, false-source, open-gap); status is open, wip (elimination
% underway), resolved, or historical. Machine-read by the site index;
% prints a verdict line.
\newcommand{\gapnote}[2]{\par\noindent\textbf{Verdict:} #1 (#2).\par}

\title{The Actual Exterior Collar Gives an Untwisted Boundary Component}
\author{}
\date{}
\begin{document}
\maketitle
\gapnote{scope-restriction}{open}

Schuch, Cirac, and P\'erez-Garc\'ia,
\href{https://arxiv.org/abs/1001.3807v3}{arXiv:1001.3807v3},
Theorem 6.9 and Lemma 6.14, source
\path{Papers/1001.3807/paper_v3.tex}, lines 1935--1990 and 2199--2214,
use an exterior block and virtual invariance to compare flux strings away
from their endpoints. For a noncentral insertion, identity on the observed
incident bonds alone does not supply the complementary contraction. A
common untwisted component connecting all complementary boundary endpoints
is an independent geometric requirement.

The theorem
\leanid{TNLean.PEPS.torusDualFlux_transport_eq_one_on_integerExteriorCollar}
derives identity transport on every projected safe exterior collar edge
from actual dual-path avoidance. A safe edge has two unoccupied integer
endpoints near a common occupied center. That center is a corner of one
of the two adjacent dual plaquettes. A dual path avoiding all plaquettes
with observed corners cannot cross the edge, so its signed crossing count
is zero. This proves the group-valued identity without centrality.

The theorem
\leanid{TNLean.PEPS.exists_torusFlux_untwisted_boundary_component}
assumes simple connectedness of the observation region's actual closed-cell
realization and avoidance of the computed four-corner forbidden plaquettes
by the given native path. It derives a simultaneous planar integer lift,
its connected safe exterior collar, a projected untwisted complementary
subgraph and a root. Every complementary boundary endpoint lifts into
this same collar, and projected collar walks prove the required reachability.
No component, root, gauge, transport identity, Gram identity or density
equality is supplied as a hypothesis. Both torus periods are at least three;
the group and its insertion element need not be central or finite.

This proves the boundary-component geometry for an actually avoiding path.
Constructing that path from the original native flux string still requires
the explicit computed-forbidden-set hypotheses of the detour theorem and
its lifted displacement class. Neither endpoint exclusion alone nor equal
winding is asserted to classify restricted homotopies. The annular
Wilson-loop obstruction to arbitrary endpoint-free vacuum density is
retained. Composition with noncentral contraction and native coefficient
conversion remains separate under \ghissue{8703}; no full-lattice
Hamiltonian conclusion under \ghissue{8947} is asserted.
\end{document}
